\section{Historical Evidence and Scope}
\label{sec:historical-evidence}
Two archived campaigns motivate the design without constituting production
confirmation. The original FFT study holds its scalar transform implementation
and requested frame cadence fixed while comparing immediate and incremental
execution. Its numerical checks support the tested transforms and schedules;
its timings locate residual preparation and reconstruction bursts and show why
a percentile alone can miss work occurring less than once per hundred sample
calls. Appendix~\ref{sec:method} records its host, compiler, workloads, and
within-session resampling; Appendix~\ref{sec:legacy-results} retains every
reported numerical and timing summary. Those statistical choices apply only
to that campaign.

The separate complete-pipeline prototype study compares a one-hop serial
schedule with boundary bursts and a four-hop overlapped schedule
(Appendix~\ref{sec:pipeline-evaluation}). It illustrates that lower observed
block maxima can coexist with higher aggregate cost. The one-hop design also
avoids the additional $3H$ samples of age required by that overlapped
prototype. Some comparisons change storage and arithmetic as well as work
placement, so they do not identify scheduling overhead in isolation.
The selected timing examples, complete ranges, and production regression
record remain with their original evidence in the appendix.

Neither campaign measures the current weighted scheduler against optimized
libraries, concurrent display consumption, or a running audio device. The
production comparisons in Section~\ref{sec:evaluation} must supply their own
validated source identities, workload coverage, numerical audits, and session
observations. Historical data and implementation tests cannot establish a
current backend ranking or an audio-deadline guarantee.
